\chapter*{Tóm tắt}
\addcontentsline{toc}{chapter}{Tóm tắt}

Các hệ thống điều phối xe đi chung theo yêu cầu hiện nay đều báo cho hành khách một giờ đón
và một giờ trả dự kiến, rồi sau đó tiếp tục tối ưu lại kế hoạch mỗi khi có khách mới. Hệ quả
là giờ đã báo bị dời đi nhiều lần trong cùng một chuyến, trong khi mọi thước đo chuẩn của
ngành --- tỉ lệ phục vụ, thời gian chờ, độ vòng đường --- đều không ghi nhận điều đó. Một
hành khách được báo đúng một lần và một hành khách bị đổi giờ mười hai lần là giống hệt nhau
trong mọi hàm mục tiêu của dòng văn liệu hiện hành, miễn là cả hai vẫn nằm trong cửa sổ thời
gian hợp lệ.

Khóa luận này đặt câu hỏi ngược lại: \emph{một hãng gọi xe chung có thể đưa ra lời hứa về giờ
giấc mà hãng thực sự giữ được hay không, và cái đó đắt bao nhiêu?} Để trả lời, chúng tôi xây
dựng cơ chế \textbf{lời hứa có biên được chứng nhận}. Mỗi lần hệ thống công bố lại một lời
hứa, độ dời được ghi vào một sổ cái riêng của từng hành khách, cộng dồn và chỉ tăng, rồi đối
chiếu với một trần $\beta$; phương án nào làm sổ cái vượt trần sẽ bị loại khỏi tập lựa chọn
ngay tại thời điểm quyết định. Khác biệt cốt lõi so với văn liệu là ở chỗ độ dời được phân rã
theo nguyên nhân: bằng cách chiếu lại chính kế hoạch cũ từ vị trí thật hiện tại của xe, hệ
thống tách được phần do nhà điều hành tự gây ra khỏi phần do thế giới vật lý gây ra, và chỉ
tính vào sổ cái phần thứ nhất.

Phần lý thuyết chứng minh ba định lý về hai cách tính sổ cái --- tính theo hành vi của nhà
điều hành, hay tính theo kết quả mà khách nhìn thấy. Cách thứ nhất luôn thỏa được; cách thứ
hai có thể khiến hệ thống không còn bất kỳ quyết định hợp lệ nào, và không bao giờ phục vụ
được nhiều hơn cách thứ nhất với bất kỳ trần $\beta$ nào. Phần thực nghiệm gồm 558 lần chạy
mô phỏng trên dữ liệu chuyến đi công khai của Manhattan, sử dụng nền mô phỏng FleetPy và một
panel 20 ô đã niêm phong. Kết quả chính cho thấy $89{,}06\%$ lượng dịch chuyển giờ hẹn mà hành
khách nhìn thấy trong chế độ không cam kết là do chính bộ điều phối gây ra; một trần 30 giây
cắt con số đó xuống $1{,}37\%$ với cái giá $7{,}13$ điểm phần trăm tỉ lệ phục vụ; và dưới dữ
liệu tắc đường đo thật của Manhattan, cách tính theo kết quả làm hệ thống tê liệt trên 19 trên
20 ô, trong khi cách tính theo hành vi không hỏng một ô nào.

\vspace{6pt}
\noindent\textbf{Từ khóa:} điều phối xe đi chung, lời hứa có biên, độ ổn định lịch trình,
quy trách nguyên nhân, nghĩa vụ phương tiện, đánh đổi phục vụ--ổn định.

\clearpage

\chapter*{Abstract}
\addcontentsline{toc}{chapter}{Abstract}

Modern on-demand ride-pooling dispatchers announce an estimated pickup and drop-off time to
every rider, then keep re-optimising the plan as new requests arrive. As a result, the
announced times are revised repeatedly within a single trip, while none of the field's
standard metrics --- service rate, waiting time, detour --- record this at all. A rider told
once and a rider re-quoted twelve times are indistinguishable under every objective function
in the current literature, as long as both remain inside the admissible time window.

This thesis asks the opposite question: \emph{can a ride-pooling operator make a timing
promise it can actually keep, and what does that cost?} We build a mechanism of
\textbf{certified bounded promises}. Every time the system republishes a promise, the
displacement is recorded in a per-rider ledger that accumulates monotonically and is checked
against a cap $\beta$; any candidate plan that would push the ledger past the cap is removed
from the choice set at decision time. The core departure from the literature is that the
displacement is decomposed by cause: by re-projecting the unchanged previous plan from the
vehicle's true current position, the system separates what the operator caused from what the
physical world caused, and charges only the former.

The theoretical part proves three theorems about two ways of keeping the ledger --- charging
the operator's conduct, or charging the outcome the rider observes. The first is always
satisfiable; the second can leave the dispatcher with no admissible decision at all, and never
serves more riders than the first for any cap $\beta$. The experimental part comprises 558
simulation runs on public Manhattan trip data, using the FleetPy simulator and a sealed
20-cell panel. The headline results: $89.06\%$ of the rider-visible promise movement under the
unconstrained policy is operator-induced; a 30-second cap reduces this to $1.37\%$ at a cost of
$7.13$ percentage points of service rate; and under measured Manhattan congestion, the
outcome basis renders the system infeasible on 19 of 20 cells while the conduct basis fails on
none.

\vspace{6pt}
\noindent\textbf{Keywords:} ride-pooling dispatch, bounded promises, schedule stability, causal
attribution, obligation of means, service--stability trade-off.

\clearpage
